% JoonHo Lee (jlee296@ua.edu)
% Requires booktabs, tabularx, and array.
\begin{table}[htbp]
\caption{Program year 2026 rating form for the Small City, Large City, and County funds}
\label{tab:F-rating}
\centering\small
\renewcommand{\arraystretch}{1.2}
\begin{tabularx}{\linewidth}{@{}>{\raggedright\arraybackslash}p{.25\linewidth}>{\raggedright\arraybackslash}X>{\raggedright\arraybackslash}X@{}}
\toprule
Component & Points & Evidence assessed \\
\midrule
Needs assessment & 5 & Community-wide need and low- and moderate-income conditions. \\
\addlinespace[0.35em]
Project development & 45 & Need, alternatives, activities, and costs; includes 10 points for project maps. \\
\addlinespace[0.35em]
Impact & 50 & Actual beneficiaries, directness, and urgency. \\
\addlinespace[0.35em]
Other considerations & 30 & Operations, maintenance, participation, implementation capacity, and distress. \\
\addlinespace[0.35em]
\midrule
Nature of benefits & 130 & Subtotal of the four components above. \\
\addlinespace[0.35em]
Cost-benefit ratio & 50 & Requested dollars per actual beneficiary. \\
\addlinespace[0.35em]
Matching contribution & 20 & Local match; full score for applicants with populations of 1,000 or fewer, or projects with 100\% LMI beneficiaries. \\
\addlinespace[0.35em]
\midrule
Overall cap & 200 & Total score, subject to the form's provisions. \\
\addlinespace[0.35em]
\bottomrule
\end{tabularx}
\par\medskip
\begin{minipage}{\linewidth}\small
\textit{Note.} The 45-point Project Development category includes 10 points for maps. Nature of Benefits is the 130-point subtotal, not an additional category. The overall cap is 200. Bonus and adjustment provisions are outside this summary; completed project scores are unobserved. LMI means low and moderate income.
\end{minipage}
\end{table}
